% GENERATED FILE. Edit canonical lessons, then run npm run generate:books.
% canonical-source-hash: fnv1a64:a48f60fe137a11f1
% canonical-lessons: IT-C397-garanzia, IT-C397-garantire, IT-C397-fattura, IT-C397-preventivo, IT-C397-conveniente

\chapter{Garanzie e fatture}
\label{ch:it-garanzie-e-fatture}
\hlchaptermodality{397}

\begin{chapteropening}
\textbf{By the end of this chapter:} \emph{I can deal with what I buy: a guarantee, an invoice, a quote and whether it is good value.}

Complete the last lesson of chapter 397: I can deal with what I buy: a guarantee, an invoice, a quote and whether it is good value.
\end{chapteropening}

\section[la garanzia]{la garanzia — the guarantee, the warranty}

\label{lesson:IT-C397-garanzia}



\begin{quote}
Before the new one: say the Italian for the payment, then the Italian for the deposit.
\end{quote}

\subsection*{You'll want to know: la garanzia}
\textbf{la garanzia} — \textquotedblleft{}the guarantee, the warranty\textquotedblright{}.

\begin{quote}
Se si rompe, è ancora in garanzia.
\end{quote}

\begin{quote}
\emph{If it breaks, it is still under warranty.}
\end{quote}

\emph{In garanzia} — under warranty.

\begin{cousinweb}
\emph{garanzia} is built on \emph{garante} \textquotedblleft{}guarantor\textquotedblright{}, which came from French \emph{garant}, a word of Germanic, probably Frankish, origin; EtymDB sets it beside Old High German \emph{werento} \textquotedblleft{}guarantor\textquotedblright{} (EtymDB), but the exact form behind it is not settled.
\end{cousinweb}

\begin{grammarlens}[title={What you've built}]
One more word for telling, explaining and describing at B1.
\end{grammarlens}

\subsection*{Your turn}
\begin{itemize}
\raggedright
  \item \emph{Say it:} \emph{la garanzia}
  \item \emph{Say it:} \emph{la garanzia}, once more
  \item \emph{Say it:} a sentence of your own with \emph{la garanzia}
  \item \emph{Recall:} say the Italian for the payment, then the Italian for the deposit, then say \emph{la garanzia} again
\end{itemize}

\begin{tcolorbox}[breakable,colback=teal!4,colframe=teal!35!black,title={Before you move on}]
What does la garanzia mean? (\textquotedblleft{}The guarantee, the warranty\textquotedblright{}.) Say it once more.
\end{tcolorbox}
\section[garantire]{garantire — to guarantee}

\label{lesson:IT-C397-garantire}



\begin{quote}
Before the new one: say the Italian for the deposit, then the Italian for the guarantee.
\end{quote}

\subsection*{You'll want to know: garantire}
\textbf{garantire} — \textquotedblleft{}to guarantee\textquotedblright{}.

\begin{quote}
Ti garantisco che domani è pronto.
\end{quote}

\begin{quote}
\emph{I guarantee you it will be ready tomorrow.}
\end{quote}

Conjugated with \emph{-isc-}: \emph{garantisco}.

\begin{cousinweb}
\emph{garantire} is built on \emph{garante} \textquotedblleft{}guarantor\textquotedblright{}, from French \emph{garant}, a word of Germanic, probably Frankish, origin; EtymDB sets it beside Old High German \emph{werento} \textquotedblleft{}guarantor\textquotedblright{} (EtymDB), but the exact form behind it is not settled.
\end{cousinweb}

\begin{grammarlens}[title={What you've built}]
One more word for telling, explaining and describing at B1.
\end{grammarlens}

\subsection*{Your turn}
\begin{itemize}
\raggedright
  \item \emph{Say it:} \emph{garantire}
  \item \emph{Say it:} \emph{garantire}, once more
  \item \emph{Say it:} a sentence of your own with \emph{garantire}
  \item \emph{Recall:} say the Italian for the deposit, then the Italian for the guarantee, then say \emph{garantire} again
\end{itemize}

\begin{tcolorbox}[breakable,colback=teal!4,colframe=teal!35!black,title={Before you move on}]
What does garantire mean? (\textquotedblleft{}To guarantee\textquotedblright{}.) Say it once more.
\end{tcolorbox}
\section[la fattura]{la fattura — the invoice}

\label{lesson:IT-C397-fattura}



\begin{quote}
Before the new one: say the Italian for the guarantee, then the Italian for to guarantee.
\end{quote}

\subsection*{You'll want to know: la fattura}
\textbf{la fattura} — \textquotedblleft{}the invoice\textquotedblright{}.

\begin{quote}
Se mi fa la fattura, posso scaricare la spesa.
\end{quote}

\begin{quote}
\emph{If you give me an invoice, I can deduct the cost.}
\end{quote}

\emph{Fare la fattura} — to issue an invoice. It can also mean a spell or a curse.

\begin{cousinweb}
\emph{fattura} is inherited from Latin \emph{factūra} \textquotedblleft{}a making, a piece of work\textquotedblright{} (Lewis \& Short), from \emph{facere} \textquotedblleft{}to make\textquotedblright{}: the record of the work done.
\end{cousinweb}

\begin{grammarlens}[title={What you've built}]
One more word for telling, explaining and describing at B1.
\end{grammarlens}

\subsection*{Your turn}
\begin{itemize}
\raggedright
  \item \emph{Say it:} \emph{la fattura}
  \item \emph{Say it:} \emph{la fattura}, once more
  \item \emph{Say it:} a sentence of your own with \emph{la fattura}
  \item \emph{Recall:} say the Italian for the guarantee, then the Italian for to guarantee, then say \emph{la fattura} again
\end{itemize}

\begin{tcolorbox}[breakable,colback=teal!4,colframe=teal!35!black,title={Before you move on}]
What does la fattura mean? (\textquotedblleft{}The invoice\textquotedblright{}.) Say it once more.
\end{tcolorbox}
\section[il preventivo]{il preventivo — the quote, the estimate}

\label{lesson:IT-C397-preventivo}



\begin{quote}
Before the new one: say the Italian for to guarantee, then the Italian for the invoice.
\end{quote}

\subsection*{You'll want to know: il preventivo}
\textbf{il preventivo} — \textquotedblleft{}the quote, the estimate\textquotedblright{}.

\begin{quote}
Prima di cominciare, l'idraulico ci fa un preventivo.
\end{quote}

\begin{quote}
\emph{Before starting, the plumber gives us a quote.}
\end{quote}

\emph{Chiedere un preventivo} — to ask for a quote.

\begin{cousinweb}
\emph{preventivo} is a learned formation on Latin \emph{praevent-}, from \emph{praevenīre} \textquotedblleft{}to come before, to anticipate\textquotedblright{} (Lewis \& Short), built on \emph{venīre} \textquotedblleft{}to come\textquotedblright{}: the cost reckoned beforehand.
\end{cousinweb}

\begin{grammarlens}[title={What you've built}]
One more word for telling, explaining and describing at B1.
\end{grammarlens}

\subsection*{Your turn}
\begin{itemize}
\raggedright
  \item \emph{Say it:} \emph{il preventivo}
  \item \emph{Say it:} \emph{il preventivo}, once more
  \item \emph{Say it:} a sentence of your own with \emph{il preventivo}
  \item \emph{Recall:} say the Italian for to guarantee, then the Italian for the invoice, then say \emph{il preventivo} again
\end{itemize}

\begin{tcolorbox}[breakable,colback=teal!4,colframe=teal!35!black,title={Before you move on}]
What does il preventivo mean? (\textquotedblleft{}The quote, the estimate\textquotedblright{}.) Say it once more.
\end{tcolorbox}
\section[conveniente]{conveniente — good value, cheap; worthwhile}

\label{lesson:IT-C397-conveniente}



\begin{quote}
Before the new one: say the Italian for the invoice, then the Italian for the quote.
\end{quote}

\subsection*{You'll want to know: conveniente}
\textbf{conveniente} — \textquotedblleft{}good value, cheap; worthwhile\textquotedblright{}.

\begin{quote}
Il treno è più conveniente della macchina.
\end{quote}

\begin{quote}
\emph{The train is better value than the car.}
\end{quote}

A false friend: it does not mean convenient (that is \emph{comodo}).

\begin{cousinweb}
\emph{conveniente} comes from Latin \emph{conveniēns} \textquotedblleft{}agreeing, suitable\textquotedblright{} (EtymDB), present participle of \emph{convenīre} \textquotedblleft{}to come together, to suit\textquotedblright{}, built on \emph{venīre} \textquotedblleft{}to come\textquotedblright{}.
\end{cousinweb}

\begin{grammarlens}[title={What you've built}]
One more word for telling, explaining and describing at B1.
\end{grammarlens}

\subsection*{Your turn}
\begin{itemize}
\raggedright
  \item \emph{Say it:} \emph{conveniente}
  \item \emph{Say it:} \emph{conveniente}, once more
  \item \emph{Say it:} a sentence of your own with \emph{conveniente}
  \item \emph{Recall:} say the Italian for the invoice, then the Italian for the quote, then say \emph{conveniente} again
\end{itemize}

\begin{tcolorbox}[breakable,colback=teal!4,colframe=teal!35!black,title={Before you move on}]
What does conveniente mean? (\textquotedblleft{}Good value, cheap; worthwhile\textquotedblright{}.) Say it once more.
\end{tcolorbox}
